\section{Methodik und Experimentdesign}

\subsection{Stage A: abgeschlossener Klassifikationsvorversuch}
Stage~A verwendete 796 unter kontrollierten Bedingungen aufgenommene Bilder der Klassen \textit{glass, metal, paper} und \textit{plastic}. Mit Seed~123 wurden \SI{80}{\percent} für Training und \SI{20}{\percent} für Validierung vorgesehen (etwa 637/159 Bilder). Verglichen wurden ein kleines CNN, MobileNetV2 mit eingefrorenem Backbone, Fine-Tuning und der INT8-Export \cite{sandler2018mobilenetv2,tensorflow2024transfer}. Eingabegrößen und Vorverarbeitung waren innerhalb des jeweiligen Modellvergleichs festgelegt. Stage~A bewertet Bildklassifikation mit Accuracy; ihre Resultate sind nicht mit mAP-Werten der Detektion gleichzusetzen.

\subsection{Historische Detektionsentwicklung}
\label{sec:method-gigo}
Für die Objekterkennung werden Bounding Boxes benötigt. Ein früher Auto-Labeling-Versuch mit Canny-Kanten und Otsu-Schwellenwert erfasste häufig kontrastreiche Tischbereiche statt der Gegenstände. Das zugehörige EXP-005 war deshalb kein belastbarer Nachweis der Generalisierung. Danach wurden unter anderem TACO, Roboflow Trash Detection und ein mithilfe von SAM nachannotierter TrashNet-Bestand eingesetzt \cite{proenca2020taco,roboflow2023trash,yang2016classification,kirillov2023segment}.

Die YOLO11n-Läufe EXP-013 bis EXP-017 sind historische Experimente. EXP-013 nutzte TACO und TrashNet, EXP-014 zusätzlich Roboflow, EXP-015 TACO, TrashNet und WaRP, EXP-016 nur WaRP und EXP-017 alle vier Quellen \cite{warp2024dataset,ultralyticsyolov8}. Sie liefen jeweils 120 Epochen mit $640\times640$ Pixeln. Da sich Zusammensetzung und resultierende Validierungsbestände änderten, werden ihre mAP-Werte deskriptiv berichtet, aber nicht als kontrollierte Ablation oder als Resultat des aktuellen Datensatzes interpretiert.

\subsection{Aktueller balancierter Fünfklassen-Datensatz}
Der aktuelle Bestand \texttt{merged\_mira\_balanced} enthält 6898 Bilder: 5108 Training, 415 Validierung und 1375 Test. Alle Quellklassen werden programmatisch auf \textit{glass, metal, paper, plastic} und \textit{trash} abgebildet. Die Quellen und Splitregeln sind:
\begin{itemize}
    \item TACO: deterministische Aufteilung 70/15/15 mit Seed~42 \cite{proenca2020taco};
    \item Roboflow Trash Detection: Übernahme der veröffentlichten Train/Valid/Test-Zuordnung \cite{roboflow2023trash};
    \item SAM-annotiertes TrashNet: vorhandene Train/Val-Zuordnung \cite{yang2016classification,kirillov2023segment};
    \item dmedhi Garbage Image Classification/Detection: Übernahme der veröffentlichten Train/Validation-Zuordnung \cite{dmedhiGarbageDataset}.
\end{itemize}

Das Training wird aus den jeweiligen Trainingsanteilen gebildet. Ein deterministischer Auswahlalgorithmus mit Seed~42 nähert die Klassen nach Anzahl der Bounding Boxes bis zu einem Ziel von höchstens 1800 Boxen je Klasse an. Da Bilder mehrere Objekte enthalten können und Dubletten anschließend entfernt werden, ist das Ergebnis annähernd, nicht exakt, ausgeglichen. Der Validierungssatz besteht ausschließlich aus dem zurückgehaltenen TrashNet-Val-Anteil und bildet damit das geplante Tischszenario ab. Der Testsatz kombiniert dmedhi-Validation sowie TACO- und Roboflow-Val/Test-Anteile. SHA-256-Prüfsummen verhindern identische Bilddateien über Splitgrenzen hinweg; sie schließen inhaltlich ähnliche Aufnahmen jedoch nicht sicher aus. Ein zeilenweises Manifest speichert für jedes Bild Zielsplit, Quelle, ursprünglichen Split, Quell-ID, Annotationen und Prüfsumme.

\subsection{Geplanter kontrollierter Vergleich}
Auf dem aktuellen Datensatz wird YOLO11n trainiert und anschließend in FP32- und INT8-Form evaluiert. Der Lauf ist zum Zeitpunkt dieser Fassung noch nicht abgeschlossen. Für beide Varianten bleiben Testdaten, Eingabegröße, Konfidenzregeln und Messhardware gleich. Berichtet werden mAP50--95 als Hauptmetrik, mAP50 und klassenweise AP sowie Precision, Recall und $F_1$ bei dokumentierter Konfidenzschwelle. Zusätzlich werden Dateigröße und CPU-Latenz nach Aufwärmläufen über dieselbe Bildmenge gemessen. Dadurch lässt sich der Quantisierungseffekt als Differenz beziehungsweise Verhältnis angeben, ohne Hardware- und Datensatzeffekte zu vermischen.

Der Testsatz bleibt bis zur endgültigen Modellwahl unangetastet; Anpassungen erfolgen anhand des Validierungssatzes. Ein späterer Raspberry-Pi-Benchmark muss dasselbe exportierte Modell, dieselbe Eingabegröße und ein festes Messprotokoll verwenden. Kamera, Tracking und Roboterarm werden danach separat untersucht, damit mechanische und softwareseitige Fehler nicht vermischt werden.
